\scanpage{409}
\mspage{323}
% archive sheet number: 408

% 1–6
Ainsi, on a construit les deux objets et les quatre flèches voulues, de façon à satisfaire aux \unsure{conditions|compatibilités|contraintes} $1_\alpha,2_{\unsure{\alpha}}$ de la page 318. Les contraintes $3_\alpha(\beta),4_\alpha(\beta)$ sont-elles alors \unsure{satisfaites} automatiquement ?
% 7–9
\diagram{\begin{tikzcd}[ampersand replacement=\&,column sep=large]
\varphi(i_\beta)\arrow[r,"u_\beta"]\arrow[d] \& \unsure{j|J}_\beta\arrow[d]\\
\varphi(i_{\alpha_0})\arrow[r,"u_{\alpha_0}"]\arrow[d,"{\varphi(\tau_{\alpha,\alpha_0})}"] \& \deleted{\unsure{\varphi(}}j_{\alpha_0}\deleted{\unsure{)}}\arrow[d,"{\tau'_{\alpha,\alpha_0}}"]\\
\varphi(i_\alpha)\arrow[r,"u_\alpha"] \& j_\alpha
\end{tikzcd}}

% 7–8, right
\margin{$\beta<\alpha$ i.e.\\ $\beta\leq\alpha_0$}

% 10–14
Le carré supérieur est commutatif par hyp. de récurrence, pour \deleted{\unsure{voir}} que le carré total \inserted{(composé)} soit commutatif, il suffit \deleted{de le \unsure{v}} que le carré inférieur le soit. On a donc une double-flèche
% 15
\[
\deleted{\unsure{\text{dans}}}\quad(\tau'_{\alpha,\alpha_0}u_{\alpha_0},u_\alpha\varphi(\tau_{\alpha,\alpha_0})):\varphi(i_{\alpha_0})\mathrel{\substack{\unsure{\hookrightarrow|\leftrightarrow|\rightarrow}\\[-.55ex]\rightarrow}}j_\alpha,
\]
% 16–19
elle n'est peut-être pas triviale mais on sait, $J'$ étant \unsure{filtrant|filtrante}, qu'on peut l'égaliser par une flèche $j_\alpha\to\overrightarrow{j}_\alpha$ dans $J'$. Donc quitte à remplacer \deleted{\unsure{$j$}}$j_\alpha$ par
